\subsection{Median on the Plane}

\subsubsection{Описание решения}

Исходя из условия задачи, для любой точки \( p_1\in P \) можно выбрать \( p_2\in P \) так, что прямая \( p_1p_2 \) разделит точки из \( P\backslash\{p_1,p_2\} \) на две равновеликие части. Это верно, потому что \( \abs{P}\divisible 2 \) и любые три точки неколлинеарны по условию.

Для определённости выберем \( p_1 \) с наименьшей абсциссой. Таких точек не более двух --- выберем любую из них и исключим из исходного множества.

Далее отсортируем точки \( p'\in P\backslash\{p_1\} \) по угловым коэффициентам прямых \( p_1p' \). Здесь может возникнуть проблема: если в множестве \( P\backslash\{p_1\} \) всё-таки есть точка \( \tilde{p_1} \) с такой же абсциссой, как у \( p_1 \), то угловой коэффициент прямой \( p_1\tilde{p_1} \) стремится к \( \pm\infty \). Аккуратно обработаем этот случай так, чтобы не возникло деления на ноль и при этом \( \tilde{p_1} \) оказалась крайним элементом в отсортированном массиве.

Ответ к задаче --- медиана отсортированного массива \( P\backslash\{p_1\} \).

\textit{Замечание.} Стоит отметить, что решение можно оптимизировать: достаточно найти \( (\frac{\abs{P}}{2}+1) \)-ую порядковую статистику.

\subsubsection{Асимптотическая сложность}

\textit{Временная сложность решения} --- линейно-логарифмическая \( \mathcal{O}(N\cdot \log N) \). Поиск точки \( p_1 \) с минимальной абсциссой линеен, операция удаления точки из массива тоже линейна. Остаётся сортировка массива, которая происходит за линейно-логарифмическое время. 

\textit{Пространственная сложность решения} --- линейная \( \mathcal{O}(N) \). В памяти хранится массив из всех введённых точек, который в дальнейшем обрабатывается с использованием константного числа переменных.

\subsubsection{Листинг кода}

\begin{lstlisting}
#include <algorithm>
#include <cstddef>
#include <cstdint>
#include <iostream>
#include <iterator>
#include <stdexcept>
#include <vector>

class Fraction {
public:
  Fraction(int num, int denum) : numerator_(num), denumerator_(denum) {
  }

  bool operator<(const Fraction& other) const {
    return this->numerator_ * other.denumerator_ < this->denumerator_ * other.numerator_;
  }

private:
  int64_t numerator_ = 1;
  int64_t denumerator_ = 1;
};

class Point {
public:
  int x = 0;
  int y = 0;
  std::size_t idx = 0;

  friend Fraction operator-(const Point& lpt, const Point& rpt) {
    return {rpt.y - lpt.y, rpt.x - lpt.x};
  }
};

namespace {
Point FindLeftmostPoint(const std::vector<Point>& array) {
  if (array.empty()) {
    throw std::invalid_argument("array is empty");
  }
  Point left_point = array[0];
  for (std::size_t i = 1; i < array.size(); i++) {
    if (array[i].x < left_point.x) {
      left_point = array[i];
    }
  }
  return left_point;
}
}  // namespace

int main() {
  std::size_t n_points = 0;
  std::cin >> n_points;

  std::vector<Point> points;
  points.reserve(n_points);
  for (std::size_t i = 0; i < n_points; i++) {
    Point curr_point;
    std::cin >> curr_point.x >> curr_point.y;
    curr_point.idx = i;
    points.push_back(curr_point);
  }

  if (points.size() == 2) {
    std::cout << 1 << " " << 2;
    return 0;
  }

  Point left_point = {};
  try {
    left_point = FindLeftmostPoint(points);
  } catch (const std::invalid_argument& e) {
    std::cout << "Problem constraints are not met.";
    return 1;
  }
  points.erase(std::next(points.begin(), static_cast<std::ptrdiff_t>(left_point.idx)));

  std::sort(points.begin(), points.end(), [left_point](const Point& lpt, const Point& rpt) {
    return (left_point - lpt) < (left_point - rpt);
  });
  std::cout << left_point.idx + 1 << " " << points[points.size() / 2].idx + 1;
  return 0;
}
\end{lstlisting}